\documentclass[11pt,a4paper]{article}
\usepackage[T1]{fontenc}
\usepackage{lmodern,amsmath,amssymb,amsthm,mathtools,booktabs,microtype}
\usepackage[margin=27mm]{geometry}
\usepackage[colorlinks=true,linkcolor=blue,citecolor=blue,urlcolor=blue]{hyperref}
\newtheorem{theorem}{Theorem}[section]
\newtheorem{lemma}[theorem]{Lemma}
\newtheorem{proposition}[theorem]{Proposition}
\newcommand{\Q}{\mathbb Q}\newcommand{\Z}{\mathbb Z}
\DeclareMathOperator{\Tr}{Tr}\DeclareMathOperator{\Nm}{N}
\DeclareMathOperator{\rad}{rad}
\setlength{\emergencystretch}{2em}
\title{A Quartic in Eighty-Nine Variables under the Rational $abc$ Conjecture}
\author{}\date{10 November 2026}
\begin{document}\maketitle
\begin{abstract}
Assuming the rational $abc$ conjecture and the specified elliptic-group,
pole-realization and simultaneous prime-pattern inputs, we construct a
family of integer quartics in eighty-nine total variables. Every
sufficiently large fixed odd input shift gives rational input projection
exactly $\Z$. The quartic is a sum of seventy-seven primitive quadratic
squares and has exactly $936$ monomials. Its coefficient height, for a
shift $B\ge3$, is $\max(4153456,B^2+797681)$. A common elliptic chart in four field coordinates stores both reciprocal
formal parameters. Eight rational contacts symmetric about one permit a short quadratic evaluation and normalization.
A height estimate for a sum of fourth powers bounds the input together
with one elliptic index. All six local tests have order one; only plain
elliptic formal parameters are used. Reciprocal factor equations
avoid separately stored numerator fourth powers, and their zero stratum
is excluded by the point-pair equations. Three shifted contact outputs
and three inverse weights suffice for each detector, with an exact
rescaling that preserves its pole input. The sufficient shift is ineffective.
\end{abstract}

\section{Statement and arithmetic assumptions}

For a reduced fraction $u/v$, with $v>0$, put
$h(u/v)=\log\max(|u|,v)$. Write $h(1:u:v)$ for the logarithmic
projective height of the indicated rational point. The coefficient height
of an integer polynomial is the maximum absolute value of its coefficients.
All variable counts include the distinguished input.

\begin{theorem}\label{thm:main}
Under the assumptions below, there is a computable family
\[
 P_B\in\Z[X,Y_1,\ldots,Y_{88}],\qquad B\ge3\text{ odd},
\]
such that every member has degree exactly four, content one, exactly
$936$ monomials, and coefficient height
\[
 H(P_B)=\max(4153456,B^2+797681).
\]
There is a constant $B_*$ such that, for every odd $B\ge B_*$,
\[
 x\in\Z\quad\Longleftrightarrow\quad
 \exists\boldsymbol y\in\Q^{88},\quad P_B(x,\boldsymbol y)=0.
\]
The integer $B$ is a fixed coefficient parameter. No numerical sufficient
value of $B_*$ is determined. In particular, the coefficient fixture
$B=3$ is not asserted to define the integers.
\end{theorem}

Let $K=\Q(\rho)$, where $\rho^2=2$, let $\sigma$ be conjugation, and put
\[
 d=3-\rho,\quad D=d^2=11-6\rho,\qquad E:y^2=x^3-Dx,
\]
\[
 R_0=(1+\rho,2\rho),\qquad
 P=2R_0=(x_P,y_P)
   =\left(\frac{57-28\rho}{8},\frac{-660+417\rho}{32}\right).
\]
The arithmetic assumptions have the following precise scope.
\begin{enumerate}
\item The rational $abc$ conjecture: for every $\epsilon>0$ there is
$K_\epsilon\ge1$ such that $c\le K_\epsilon\rad(abc)^{1+\epsilon}$
for pairwise coprime positive integers $a+b=c$.
\item The elliptic source gives $2E(K)=\Z P$, with $P$ of infinite
order, and the index-height comparison stated in~\cite{Dependencies}.
The square-abscissa characterization proved in~\cite{Squares} then
applies to all nonzero multiples of $P$.
\item For a nonzero rational $b$, set $B_0=17b$. The pole equations
\begin{equation}\label{eq:originalpole}
 \begin{split}
 p_0^2&=B_0(1-e_0)(1+3e_0),\\
 B_0q_0^2&=(s_0^2-B_0e_0(B_0-b^2))
                  (s_0^2+3B_0e_0(B_0-b^2))
 \end{split}
\end{equation}
exclude negative odd valuations of $b$ outside $\{2,3,17\}$.
For inputs satisfying
\begin{equation}\label{eq:prepared}
 \begin{gathered}
 0<b<17,\qquad b\text{ has square reduced denominator},\\
 b\in\Q_2^{\times2}\cap\Q_3^{\times2}\cap\Q_{17}^{\times2},
 \qquad v_2(b)\ge3,\quad v_3(b)\ge1,\quad v_{17}(b)\ge2,
 \end{gathered}
\end{equation}
the retained restricted-converse and descent hypotheses give a witness
with $p_0q_0s_0\ne0$. The hypotheses and their nonzero-witness consequence
are specified in~\cite{Dependencies,Source27,PoleDensity}.
\item The simultaneous prime-pattern input has exactly the quantified
form in~\cite{Dependencies,Selection}: arbitrary fixed rational contact
coefficients after denominator clearing, the prescribed Gaussian
congruences and conjugate reduction orders, fixed avoidance and
divisibility, and compatible real approximation. Its analytic source is
Sun's specified Theorem~4.1 and the estimates pinned in those papers.
All scales and finite preparation data are fixed before prime selection.
\end{enumerate}
The exponent-nineteen contact-height input is not used here.

The one-point case of the proved aggregate trace-height theorem
\cite{Trace}, obtained from von Buhren's family height theorem, gives
constants $A>0$, $J\ge0$ such that
\begin{equation}\label{eq:traceheight}
 n^2\le A h(T(n))+J,\qquad
 T(n)=\Tr x([n]P),\quad n\in\Z\setminus\{0\}.
\end{equation}
The additive constant is ineffective. Each abscissa is a positive square
at both real embeddings and exceeds $d$ there, so $T(n)>6$.

\section{Two height consequences of $abc$}

Set
\[
 \Theta_8=\{1\pm1,1\pm2,1\pm3,1\pm4\}
          =\{-3,-2,-1,0,2,3,4,5\}.
\]
For $t\notin\Theta_8$, let $M(t)$ be the product of primes $p>76176$
for which $v_p(t-\theta)>0$ is odd for some $\theta\in\Theta_8$.

\begin{lemma}\label{lem:contact}
Assume rational $abc$. There is a fixed $C_0\ge0$ such that
\[
                         h(t)\le\frac13\log M(t)+C_0.
\]
\end{lemma}
\begin{proof}
For a primitive pair $(u,v)$, $v>0$, use the squarefree binary form
$G(u,v)=\prod_{\theta\in\Theta_8}(u-\theta v)$ of degree eight.
Granville's consequence of $abc$~\cite[Proposition~1]{Granville},
with $\epsilon=1/2$, gives
\[
 \log\rad(G(u,v))\ge\frac{11}{2}\log\max(|u|,v)-C.
\]
For a nonzero integer $m$, let $\operatorname{sf}(m)$ be the product
of primes having odd valuation in $m$. Prime by prime,
$\rad(m)^2\le |m|\operatorname{sf}(m)$.
Also $|G(u,v)|\le C'\max(|u|,v)^8$.
Outside the fixed cutoff, no prime divides two distinct factors of
$G$, and a prime dividing $v$ divides none of them. Its squarefree
part outside the cutoff is therefore precisely $M(u/v)$.
The omitted squarefree factors are bounded by a fixed product of primes.
Consequently
\[
 \log M(u/v)\ge\left(2\cdot\frac{11}{2}-8\right)h(u/v)-C''
                 =3h(u/v)-C'',
\]
as required.
\end{proof}

\begin{lemma}\label{lem:pair}
Assume rational $abc$. There is a fixed $C_1\ge0$ such that
\[
                    h(1:u:v)\le\frac43h(u^4+v^4)+C_1
                    \qquad(u,v\in\Q).
\]
One may take $C_1=\frac7{12}\log2+\frac54\log K_{1/15}$.
\end{lemma}
\begin{proof}
The zero-coordinate cases follow from $h(w^4)=4h(w)$.
Otherwise write $u=gA/D_0$, $v=gB/D_0$, where
$\gcd(A,B)=\gcd(g,D_0)=1$, $g,D_0>0$. Put
$H=\max(|A|,|B|)$ and $N=A^4+B^4$, so $H^4\le N\le2H^4$.
Let $k=\gcd(N,D_0^4)$, $m_0=N/k$, $n_0=D_0^4/k$.
Then $u^4+v^4=g^4m_0/n_0$ is reduced. Write
$M=\max(g^4m_0,n_0)$.

Every prime dividing $N$ either divides $D_0$ or survives in $m_0$,
so $\rad(N)\le D_0m_0$. Apply $abc$ with $\epsilon=1/15$ to
$A^4+B^4=N$. The three positive terms are pairwise coprime, and
\[
 H^{7/4}\le K_{1/15}^{15/16}D_0m_0.
\]
Using $D_0^4=n_0N/m_0\le2n_0H^4/m_0$ gives
\[
 H^{3/4}\le2^{1/4}K_{1/15}^{15/16}m_0^{3/4}n_0^{1/4}.
\]
Moreover $\max(gH,D_0)\le2^{1/4}H(M/m_0)^{1/4}$.
Taking logarithms and combining yields
\[
 \begin{split}
 h(1:u:v)&=\log\max(gH,D_0)\\
 &\le C_1+\frac34\log m_0+\frac13\log n_0+\frac14\log M
 \le C_1+\frac43\log M.
 \end{split}
\]
This includes common numerator factors and all denominator cancellation.
\end{proof}

\section{Quadratic point equations}

Put $\beta=2X^2+B$. If $X=a/b$ is reduced, integrality of $\beta$
implies $b^2\mid2$, hence $b=1$. Thus $\beta\in\Z$ if and only if
$X\in\Z$; for integer input $\beta$ is odd, and always $\beta\ge B$.

All equations in this section are over $K$ and mean their two rational
component equations. Products of stored quantities have degree at most two.
For the root point retain $h,j,a\in K$ and impose
\begin{equation}\label{eq:root}
                       a=h^2,\qquad j^2=a^2-D.
\end{equation}
Then $A=(a,hj)=[n]P$ for some nonzero integer $n$ by the
square-abscissa theorem~\cite{Squares}. \subsection{The common reference point}

Write $a_P=x_P$, $b_P=y_P$, and set
\[
 A_P=a_P^2+D,\qquad B_P=a_P^2-D,\qquad
 C_P=D^2+6a_P^2D-3a_P^4.
\]
The constants $a_P,b_P,D,A_P,C_P$ are nonzero. The first four
nonvanishing assertions follow from the displayed coordinates and
$D=(3-\rho)^2>0$ at both real embeddings; exact arithmetic gives
\[
 C_P=-\frac{94897379}{4096}+\frac{4195821}{256}\rho,
 \qquad \Nm(C_P)=-\frac{8203357319351}{16777216}\ne0.
\]
Retain four field coordinates $S,T_H,f,F_0$. Put
\begin{align*}
 L_f&=2b_PD(a_PT_H-b_P),\\
 Q_f&=a_PA_PST_H+b_PA_PS-b_PB_PT_H+4a_P^2D,\\
 L_g&=2b_PD(a_PS+b_P),\\
 Q_g&=-a_PA_PST_H-b_PB_PS+b_PA_PT_H-4a_P^2D.
\end{align*}
Impose the three quadratic equations
\begin{equation}\label{eq:H}
 \begin{split}
 L_ff&=Q_f,\\
 L_gF_0&=Q_g,\\
 D^2fF_0+a_PD(f+F_0)+a_PST_H+b_P(S-T_H)+D&=0.
 \end{split}
\end{equation}

\begin{lemma}\label{lem:commonchart}
These equations describe exactly the points $H\in E(K)$ for which
$H$ and $H-P$ are finite with nonzero abscissas, through
\[
 S=\frac{y(H)}{x(H)},\qquad T_H=\frac{y(H-P)}{x(H-P)},\qquad
 f=\frac1{x(H)},\qquad F_0=\frac1{x(H-P)}.
\]
In particular $fF_0\ne0$ at every solution. The statement includes
all finite exceptional cases of the displayed equations.
\end{lemma}
\begin{proof}
For the algebraic calculation abbreviate $a=a_P$, $b=b_P$, $T=T_H$
and $g=F_0$. Thus $b^2=a(a^2-D)$. Define
\[
 \Phi=aS^2T^2-(a^2-D)(S^2+T^2)+2(a^2+D)ST+4aD
\]
and the polynomial expressions
\begin{align*}
 2bDf_0&=-(a^2+D)S+(a^2-D)T-bS^2-aS^2T,\\
 2bDg_0&=(a^2+D)T-(a^2-D)S-bT^2+aST^2.
\end{align*}
Direct expansion gives
\begin{align}
 Df_0^2+S^2f_0-1
   &=\frac{a^2S^2-b^2}{4ab^2D}\Phi,\label{eq:chartcurve}\\
 Dg_0^2+T^2g_0-1
   &=\frac{a^2T^2-b^2}{4ab^2D}\Phi,\nonumber\\
 L_ff_0-Q_f&=-a\Phi,\qquad L_gg_0-Q_g=a\Phi.\nonumber
\end{align}
Let $r_f=L_ff-Q_f$, $r_g=L_gg-Q_g$, and
$R=D^2fg+aD(f+g)+aST+b(S-T)+D$. The further exact identity
\begin{equation}\label{eq:chartmaster}
\begin{split}
 L_fL_gR={}&D^2r_fr_g+(aDL_g+D^2Q_g)r_f\\
          &+(aDL_f+D^2Q_f)r_g-aD^2C_P\Phi
\end{split}
\end{equation}
shows that every solution has $\Phi=0$, including the cases in which
one leading coefficient vanishes. Also
\[
 (2bD)^2R(f_0,g_0)=
 D^2\bigl((a^2-D)-aST-b(S-T)\bigr)\Phi.
\]
If $L_fL_g\ne0$, these identities give $f=f_0,g=g_0$.
If $L_f=0$, then $T=b/a$ and $\Phi=0$ forces
\[
 S=\frac{a^4-6Da^2+D^2}{2b(a^2+D)},\qquad
 L_g=-\frac{aDC_P}{a^2+D}\ne0.
\]
Thus $g=g_0=1/a$. The coefficient of $f$ in $R$ is
$D(a^2+D)/a\ne0$, so $f=f_0$. The case $L_g=0$ is symmetric.
In particular both leading coefficients cannot vanish simultaneously;
at $S=-b/a,T=b/a$ one has $Q_f=C_P\ne0$.

Equation~\eqref{eq:chartcurve} now proves $fg\ne0$ and gives the
points $H=(f^{-1},S/f)$ and $H'=(g^{-1},T/g)$ on $E$.
Two more identities prove the translation:
\[
 a(f_0T+g_0S)+b(g_0-f_0)=S+T,
\]
\[
 f_0g_0(aS+b)^2-a(1-af_0)^2
 =-\frac{(aS+b)^2\bigl(aST+b(S-T)+D\bigr)}{4b^2D^2}\Phi.
\]
If $x(H)\ne a$, the second identity gives the abscissa of
$H-P$ from the chord slope $(x(H)S+b)/(x(H)-a)$.
The first states that $H,-H',-P$ are collinear, and therefore gives
the required ordinate. If $x(H)=a$, then $S=\pm b/a$.
At $S=b/a$ the formula gives $f=-a/D\ne1/a$, which is impossible.
At $S=-b/a$ it gives $H=-P$ and
\[
 T=-\frac{a^4-6Da^2+D^2}{2b(a^2+D)},\qquad
 g=\frac{4b^2}{(a^2+D)^2},
\]
which are the coordinates of $H'=-2P$ in the stated form.
Thus $H'=H-P$ in all cases.

Conversely, the chord-and-tangent formulas for any allowed $H$ give
$\Phi=0$ and the expressions $f_0,g_0$ above; the same explicit
doubling computation covers $H=-P$. The displayed identities then
prove all three equations. This proves both directions.
\end{proof}

Write $H=[m]R_0+T_0$, with $T_0\in E[2](K)$, and put $h_*=m/2$.
This notation is an elliptic index; it does not define multiplication
of a point by a nonintegral rational number.


\subsection{A pair of points with fixed difference}

Retain $V,Q,W_A,I_A,K_A\in K$ satisfying
\begin{equation}\label{eq:pairchart}
 \begin{split}
 hI_A&=(j+DW_A)V+aQ,\\
 hK_A&=-(j+DW_A)Q-aV,\\
 VQ&=aW_A,\\
 V^2+Q^2+2jW_A+DW_A^2&=1.
 \end{split}
\end{equation}

\begin{lemma}\label{lem:pairchart}
Equations~\eqref{eq:root} and~\eqref{eq:pairchart} describe points
$U,V_1\in\Z P$ with $U-V_1=A$, allowing either numerator point
to be the identity. Both $I_A$ and $K_A$ are nonzero. On the
nonidentity branches,
\[
 U=\left(\frac a{V^2},\frac{aI_A}{V^3}\right),\qquad
 V_1=\left(\frac a{Q^2},\frac{aK_A}{Q^3}\right),
 \qquad z(U)=-\frac V{I_A},\quad z(V_1)=-\frac Q{K_A},
\]
where $z=-x/y$. In all branches the quantities $R_U,R_1\in K$
defined by
\begin{equation}\label{eq:recipsquares}
             I_A^2=a-DaR_U,\qquad K_A^2=a-DaR_1
\end{equation}
are $f(U)^2,f(V_1)^2$, with $f=1/x$ extended by zero at the identity.
\end{lemma}
\begin{proof}
The square chart gives $haj\ne0$.
Squaring the first two equations and using the last two gives
$aI_A^2=a^2-DV^4$ and $aK_A^2=a^2-DQ^4$.
For $VQ\ne0$ these are exactly the two curve equations. Their
abscissas are squares, so the points lie in $\Z P$.
The line with slope $(1/W_A-j)/h$ through $U$ also passes through
$-V_1$ and $-A$. Substitution gives its three abscissas
$a/V^2,a/Q^2,a$, counted with multiplicity. The chord-and-tangent
law therefore gives $U-V_1=A$, including the tangent case.
Nontorsion of the nonidentity points implies $I_AK_A\ne0$.

If $V=0$, then $W_A=0$, $Q^2=1$, $I_A=hQ$ and
$K_A=-jQ/h$. The interpretation is $U=O,V_1=-A$.
If $Q=0$, then $W_A=0$, $V^2=1$, $I_A=jV/h$ and
$K_A=-hV$, giving $U=A,V_1=O$. The two zero cases cannot occur
together. The asserted nonvanishing and reciprocal-square identities
also hold there.

Conversely, let $U,V_1\in\Z P$ be nonidentity points with $U-V_1=A$.
Their abscissas are squares. The line through $U,-V_1,-A$ has equation
$y=\mu x+c$, with $c=-h(\mu h+j)$. Put $R=\mu h+j$.
Vieta's identities give
$x(U)x(V_1)=R^2$ and $x(U)+x(V_1)=\mu^2-a$; in particular $R\ne0$.
Choose square roots $u_0^2=x(U)$, $v_0^2=x(V_1)$ with $u_0v_0=R$,
and put
\[
 V=h/u_0,\qquad Q=h/v_0,\qquad W_A=1/R,\qquad
 I_A=y(U)V^3/a,\qquad K_A=y(V_1)Q^3/a.
\]
The two Vieta identities give the last two chart equations, and the
line ordinates give the first two. This also covers tangent lines.
The identity branches are supplied by the explicit formulas above.
\end{proof}

\subsection{The two formal-parameter ratios}

Retain $L,L_0\in K$ and impose the quadratic equations
\begin{equation}\label{eq:ratios}
                         I_AL=VS,\qquad K_AL_0=QT_H.
\end{equation}
The pair-chart lemma gives $I_AK_A\ne0$. Once $H$ and $H-P$ are
nontorsion, the common-point lemma and $z=-x/y$ identify these values as
\[
 L=\frac{z(U)}{z(H)},\qquad
 L'=\sigma L_0=\sigma\!\left(\frac{z(V_1)}{z(H-P)}\right).
\]
The equations also apply when $U$ or $V_1$ is the identity, with its
formal parameter interpreted as zero. No additional coordinate for
either numerator parameter is required.

\section{Six tests and two comparison scales}

Set
\[
 C=L-L',\qquad I=L-\beta,\qquad T=\Tr a,
 \qquad a_0=1,\quad a_1=\beta^4+T^4.
\]
The six field tests are $C,I,f(H),f(U),\sigma f(H-P),\sigma f(V_1)$.
Write $L=L_0^{\mathrm{prim}}+\rho L_1^{\mathrm{prim}}$, with rational components.
In the quadratic source retain the rational component $I_0$ of $I$, with
\[
                  I_0-L_0^{\mathrm{prim}}+2X^2+B=0.
\]
Thus $\beta=L_0^{\mathrm{prim}}-I_0$ is affine in the retained
coordinates, and $I=I_0+\rho L_1^{\mathrm{prim}}$.
The fixed shift $B$ occurs only in this one equation.
Retain the four reciprocal squares and $C^2,I^2$ as quadratic outputs.
Their squares will occur directly in quadratic factor equations. The two fields in
\eqref{eq:recipsquares} supply two reciprocal squares directly.
Three scalar outputs compute the comparison scale without taking any
root or making any division:
\[
 s_h=\beta^2+T^2,\quad p_h=\beta T,\quad
 a_1=s_h^2-2p_h^2=\beta^4+T^4.
\]

Retain a rational $\lambda$. For $a\in\{a_0,a_1\}$ put $t=a\lambda$,
$\delta=t-1$. Use three scalar contact outputs $r_0,u_0,v_0$ and
three scalar weights $s_2,s_3,s_4$. Their literal quadratic equations are
\begin{equation}\label{eq:contactcircuit}
\begin{split}
 r_0&=t^2-2t,\qquad u_0=r_0^2+2r_0,\\
 (u_0-30r_0-v_0)s_2&=(r_0+1)s_3,\\
 u_0v_0+v_0+244u_0-820r_0&=0,
\end{split}
\end{equation}
and
\begin{equation}\label{eq:reciprocalweights}
                         s_4=s_2^2,\qquad s_3^2=s_2s_4.
\end{equation}
For interpretation, define the following rational expressions without
adjoining coordinates:
\[
 r=r_0+1=\delta^2,\qquad r_2=u_0+1=r^2,\qquad \xi=v_0-29,
\]
\[
 F=P_4(r)=(r-1)(r-4)(r-9)(r-16)
          =r^4-30r^3+273r^2-820r+576
          =\prod_{\theta\in\Theta_8}(t-\theta).
\]
The contact equations become
\begin{equation}\label{eq:unshiftedcontact}
 (r_2-30r-\xi)s_2=rs_3,\qquad r_2\xi+273r_2-820r+576=0.
\end{equation}
The affine shifts in~\eqref{eq:contactcircuit} give equivalent equations
whose constant coefficients are zero.

The stored field factors are $A_1,A_2\in K$, defined by
\begin{equation}\label{eq:clearedfactors}
\begin{split}
 s_3A_1&=s_4+s_3C^2+s_2f(H)^2+(f(U)^2)^2,\\
 s_3A_2&=s_4+s_3I^2+s_2\sigma f(H-P)^2+(\sigma f(V_1)^2)^2.
\end{split}
\end{equation}
Every fourth power on the right is evaluated directly as a quadratic
polynomial in the two components of its stored square. Store a field
value $\Omega=\omega_0+\omega_1\rho$ and scalars $\eta,b$, with
\begin{equation}\label{eq:normalizedproduct}
                         (r_0+1)\Omega=A_1A_2
\end{equation}
and
\begin{equation}\label{eq:scaledtrace}
 \omega_0^2-2\omega_0\eta-2\omega_1^2=0,
 \qquad s_3b=\eta^2.
\end{equation}
These equations are quadratic in the stored coordinates. The
unscaled factors, product, and trace parameter used in the valuation
arguments will be recovered in the following lemma; they are distinct
from $A_1,A_2,\Omega,\eta$.

The normalized pole uses seven further rational coordinates
$e,u,p,v,w,p_2,w_2$ and five quadratic equations
\begin{equation}\label{eq:normalizedpole}
 \begin{gathered}
 p_2=p^2,\qquad 17bp_2=289+34e-3e^2,\\
 e(17-b)=17uv,\qquad w_2=w^2,\qquad
 17bw_2=u^2+34uv-867v^2.
 \end{gathered}
\end{equation}
For $b\ne0$ these give~\eqref{eq:originalpole} under
\[
 e_0=e/17,\quad p_0=bp,\quad s_0=bu,\quad q_0=b^2uw.
\]
Conversely a witness with $s_0\ne0$ gives the displayed normalized
coordinates by these formulas and $v=e(17-b)/(17u)$.
This explains both the pole obstruction and its use in completeness.


\begin{lemma}\label{lem:domain}
Every rational solution has $rs_2s_3s_4F\ne0$. Define
\[
 s=\frac{s_3}{s_2},\qquad q=s^{-1},\qquad
 k=qr^2,\qquad R=\frac{F}{r^5}.
\]
Then $s_j=s^j=q^{-j}$ for $j=2,3,4$, and
\[
 q=\frac{\delta^6}{F},\qquad k=\frac{\delta^{10}}F=R^{-1}.
\]
The unscaled factors $T_i=A_i/s$ have precisely the values
\begin{equation}\label{eq:packed}
\begin{split}
 T_1&=1+qC^2+q^2f(H)^2+q^4f(U)^4,\\
 T_2&=1+qI^2+q^2\sigma f(H-P)^2+q^4\sigma f(V_1)^4.
\end{split}
\end{equation}
Put
\[
 W_0=\frac r{s^2}\Omega=T_1T_2=w_0+w_1\rho,
 \qquad \tau_0=\frac r{s^2}\eta.
\]
These recovered rational values satisfy
\begin{equation}\label{eq:trace}
                  w_0^2-2w_0\tau_0-2w_1^2=0,
                  \qquad kb=\tau_0^2.
\end{equation}
In particular the pole input is exactly $b=R\tau_0^2$.
Conversely, let $r=\delta^2$ and $qF=r^3$ with $rq\ne0$.
Values satisfying the unscaled factor, product, and trace equations
then uniquely determine the stored contact, factor, product, and
trace coordinates.
\end{lemma}
\begin{proof}
If $s_2=0$, equations~\eqref{eq:reciprocalweights} give $s_3=s_4=0$.
Equations~\eqref{eq:clearedfactors} then give
$f(U)^4=\sigma f(V_1)^4=0$. Thus both numerator reciprocal squares
vanish. The identities in the pair-chart proof give $V^4=Q^4=0$,
hence $V=Q=0$. Since $a\ne0$, the equation $VQ=aW_A$ gives
$W_A=0$, contradicting the last equation of~\eqref{eq:pairchart}.
Consequently $s_2\ne0$. The two weight equations now give
$s_3s_4\ne0$ and, with $s=s_3/s_2$,
\[
                         s^2=s_2,\quad s^3=s_3,\quad s^4=s_4.
\]
Equations~\eqref{eq:unshiftedcontact} give
$\xi=r_2-(30+s)r$ and then $P_4(r)=sr^3$.
If $r=0$, then $r_2=0$ and their last equation would give $576=0$.
Equivalently, $r_0=u_0=-1$ makes the last row of
\eqref{eq:contactcircuit} equal to $576$. Thus $r\ne0$, and
$F=P_4(r)=sr^3\ne0$. All asserted contact formulas follow.
In particular $\lambda\ne0$, since zero is a contact.

Dividing~\eqref{eq:clearedfactors} by $s^4$ and putting $T_i=A_i/s$
gives~\eqref{eq:packed}. Equation~\eqref{eq:normalizedproduct}
then gives $W_0=T_1T_2$. Multiplication of all three trace coordinates
by $r/s^2$ preserves the homogeneous trace equation. Furthermore
\[
 kb=\frac{r^2}s\frac{\eta^2}{s^3}
    =\left(\frac r{s^2}\eta\right)^2=\tau_0^2.
\]
This calculation does not divide by a factor, a product component,
or a trace parameter, and therefore also covers their possible zero
values before the local argument excludes them.

Conversely, put $s=q^{-1}$, $s_j=s^j$,
$r_0=r-1$, $u_0=r^2-1$, and
$v_0=r^2-(30+s)r+29$. Set
\[
 A_i=sT_i,\qquad \Omega=\frac{s^2}rW_0,
 \qquad \eta=\frac{s^2}r\tau_0.
\]
Direct substitution gives every displayed quadratic equation, with
the same pole input $b$. The values are uniquely determined.
\end{proof}

\section{Local extraction and denominator exclusions}

At a prime $p>76176$, the curve has good reduction and $D$ is a unit.
For a nonzero finite reciprocal abscissa $f(Q)$, its valuation at
either place of $K$ above $p$ is even. If $v(x)<0$, the curve equation
gives $2v(y)=3v(x)$; if $v(x)>0$, it gives $2v(y)=v(x)$.
The unit case is immediate. Zero reciprocal abscissas are interpreted
as having infinite valuation.
Lemma~\ref{lem:commonchart} identifies both stored reciprocals at every
solution, including the exceptional leading-coefficient cases of the
quadratic chart. Their even valuations are therefore available before
any formal-parameter quotient is interpreted.

Let $p\mid M(t)$, put $\nu=v_p(R)$, and normalize valuations by
$v(p)=1$. Then $\nu$ is positive and odd, $\delta$ is a unit, and
$v(q)=-\nu$. For each factor in~\eqref{eq:packed}, the three
nonconstant finite valuations belong respectively to
\[
                     1+2\Z,\qquad 2+4\Z,\qquad4+8\Z.
\]
Thus all term valuations, including zero for the constant term, are
distinct. Each factor is nonzero and has valuation at most zero;
valuation zero implies residue one. A zero test merely removes a term.

\begin{proposition}\label{prop:local}
At every contact prime, there is one place at which all six tests have
positive valuation. The four reciprocal-abscissa tests have valuation
at least two.
\end{proposition}
\begin{proof}
The preceding argument gives $W_0\ne0$. Equation~\eqref{eq:trace} then
implies $w_0\ne0$, since $w_0=0$ would imply $w_1=0$. Hence
\[
 \tau_0=\frac{\Nm(W_0)}{2w_0}=\frac1{\Tr(W_0^{-1})}\ne0,
 \qquad b=R\tau_0^2\ne0.
\]
Put $\ell=\max_{v\mid p}v(W_0)\le0$. The trace inequality gives
$v_p(b)\le\nu+2\ell$. This valuation is odd, so the pole obstruction
forces $0<v_p(b)\le\nu+2\ell$, hence $\ell>-\nu/2$.
At a place attaining $\ell$, each factor valuation is at least $\ell$.
Each of its terms therefore has valuation greater than $-\nu/2$.
For the three columns this gives respectively
\[
 v(C),v(I)>\nu/4,\qquad
 v(f(H)),v(\sigma f(H-P))>3\nu/4,\qquad
 v(f(U)),v(\sigma f(V_1))>7\nu/8.
\]
Integer rounding and evenness give the assertion.
\end{proof}

The finite $H$ chart may contain torsion translates. Before interpreting
either ratio, exclude
\[
                  E[2](K)\ \cup\ (P+E[2](K))\ \cup\ (-P+E[2](K)).
\]
Every permitted point in this finite set has reciprocal abscissa a unit
above $p>76176$. The torsion translation formulas, together with
\[
 \Nm(x_P)=41^2/64,\quad
 \Nm(x_P-d)=17^2/64,\quad \Nm(x_P+d)=63^2/64,
\]
give this directly. If $H$ belonged to the set, Proposition~\ref{prop:local}
would force both contact radicals to equal one. Lemmas~\ref{lem:contact}
and~\ref{lem:pair}, and $h(a_1)\le h(a_1\lambda)+h(\lambda)$, would
then give $h(1:\beta:T)\le\frac83C_0+C_1$. Since $\beta\ge B$, this is
impossible for every sufficiently large fixed $B$.

Consequently $h_*\notin\{-1,0,1\}$. Both reference points are
nontorsion, so $S T_H\ne0$ and $z(H),z(H-P)$ are defined and nonzero.
Equations~\eqref{eq:ratios} now have their stated quotient interpretation.
This exclusion does not use a logarithmic ratio and therefore precedes
its use.

\section{Soundness for all sufficiently large shifts}

Write $U=[u]P$, so $V_1=[u-n]P$. At a contact place supplied by
Proposition~\ref{prop:local}, $H,U$ lie in the primary formal group
and $H-P,V_1$ in the conjugate formal group. The formal parameter satisfies
\begin{equation}\label{eq:log}
 \log_E(Q)=z(Q)\mathcal A(Q),\qquad
 \mathcal A(Q)\equiv1\pmod{p^{4v(z(Q))}}.
\end{equation}
This is the constant-term case of the proved formal-logarithm estimate
in~\cite{Logs}. The unit is defined to be one at the identity.

The integral group identities
\[
                [m]U=[2u]H,\qquad
                [m-2]V_1=[2(u-n)](H-P)
\]
kill the torsion component of $H$. Their full logarithmic ratios are
$r=u/h_*$ and $r'=(u-n)/(h_*-1)$. At least one of $h_*,h_*-1$ is a
$p$-adic unit, so at least one full ratio is integral. The units in
\eqref{eq:log} preserve the corresponding approximate valuation, and
the positive-order test $C=L-L'$ transfers integrality to the other
approximate ratio and hence to both full ratios.
Only after this step does~\eqref{eq:log} give absolute congruences
$L\equiv r$ and $L'\equiv r'\pmod p$. Therefore
\[
 n=h_*r-(h_*-1)r'\equiv L\equiv\beta\pmod p.
\]
The argument also covers identity numerator points and half-integral
$h_*$, since $p$ is odd. In particular $\beta$ is integral at $p$.

If $\beta\ne n$, let $N_\beta$ be the positive absolute reduced
numerator of $\beta-n$. Both $M(\lambda)$ and $M(a_1\lambda)$ divide
$N_\beta$. Lemma~\ref{lem:contact} gives
\[
 h(a_1)\le h(\lambda)+h(a_1\lambda)
               \le\frac23\log N_\beta+2C_0.
\]
Put $H_2=h(1:\beta:T)$. Lemma~\ref{lem:pair} implies
\[
 H_2\le\frac89\log N_\beta+C_2,\qquad C_2=\frac83C_0+C_1.
\]
For a reduced fraction $\beta=a/b$,
$|a-nb|\le(1+|n|)\max(|a|,b)$. Since $h(\beta)\le H_2$,
\[
 \log N_\beta\le H_2+\log(1+|n|),\qquad
 H_2\le8\log(1+|n|)+9C_2.
\]
The trace-height bound~\eqref{eq:traceheight} now yields
\[
 n^2\le8A\log(1+|n|)+9AC_2+J.
\]
Its right side grows logarithmically, so $|n|$ and then $H_2$ are
uniformly bounded. Every constant here is independent of $B$.
Since $\beta\ge B$, choosing one sufficiently large odd shift excludes
$\beta\ne n$. Thus $\beta=n\in\Z$, and consequently $X\in\Z$.
Together with the earlier exceptional-point exclusion this proves
soundness for every odd $B$ beyond one fixed threshold.

\section{Simultaneous completeness}

Fix $X\in\Z$ and an odd $B\ge3$. Put $n=\beta$ and $A=[n]P$;
this fixes $T$ and both positive rational scales. Choose a nonzero
reference integer divisible by $64$ and compact real and finite-place
neighborhoods for the points
\[
 H=[h]P,\quad U=[hn]P,\quad H-P=[h-1]P,\quad
 V_1=[(h-1)n]P
\]
which avoid the finitely many chart denominators. In these neighborhoods
every rational test has a fixed lower valuation bound at the finitely
many preparation primes. All these choices precede prime selection.

Choose a fixed $D_1>0$ clearing $D_1^2\theta/a$ for both scales and
all contacts, and sufficiently divisible at the preparation primes.
Apply the retained prime-pattern input to the distinct coefficients
$-D_1^2\theta/a$, with repetitions removed, and write
\[
                \lambda=\frac{u}{D_1^2w^2},\qquad L_1=D_1^2w^2.
\]
Its prime values are $u-(\theta/a)L_1$. The zero contact makes $u$
one of the selected primes. The supplied conjugate reduction-order
conditions and the compatible congruence progression in~\cite{Dependencies}
give $h$ for which $H,U$ are formal at one place above each selected
prime and $H-P,V_1$ are formal at the conjugate place. The fixed
congruence and real neighborhoods are retained. Only natural depth one
is requested at either branch.

Both full logarithmic ratios equal $n$. Hence the four reciprocal
tests have valuation at least two and $C,I$ have valuation at least
four at the chosen place. Since $v_p(q)=-1$, every nonconstant term
of both factors in~\eqref{eq:packed} has positive valuation there.
At the other place each factor is either a unit of residue one or has
negative valuation. Thus $\Tr(W_0^{-1})$ has residue one or two and is
a unit. The pole input has valuation one at every selected contact prime.

For arbitrary positive rational $a$ the exact identity
\begin{equation}\label{eq:primeproduct}
 R=\frac{a^8\prod_{\theta\in\Theta_8}(u-(\theta/a)L_1)}
             {L_1^8(a\lambda-1)^{10}}
\end{equation}
has only square factors apart from the selected prime numerators.
In particular no square condition on $a$ is needed. At every other
prime $v_p(R)$ is even, and the reciprocal trace changes valuations
by an even number. These facts give a square reduced denominator for $b$.

At a preparation prime, increasing the fixed denominator divisibility
makes $v_p(t)\ll0$. Then
\[
 v_p(q)=v_p(R)=-2v_p(t),\qquad
 R=\delta^{-2}\prod_{\theta\in\Theta_8}
                   \left(1+\frac{1-\theta}{\delta}\right).
\]
The product is arbitrarily close to one and is a local square at
$2,3,17$. All weighted test terms are small, so $\Tr(W_0^{-1})$ is
arbitrarily close to two. Sufficient divisibility therefore gives the
three local squareclasses and valuation depths in~\eqref{eq:prepared}.

On the fixed compact real neighborhoods, as the lower bound on $\lambda$
tends to infinity, $R\sim q\sim t^{-2}$, $W_0\to1$, and
$b\sim1/(4t^2)$. Both real embeddings of every factor are positive
for sufficiently large $\lambda$. Hence $0<b<17$ and every required
denominator is nonzero. These bounds are uniform on the chosen compact
sets, so they can be fixed before applying the prime-pattern theorem.
Each prepared $b$ has a nonzero pole witness, which converts to
\eqref{eq:normalizedpole}. The converse in Lemma~\ref{lem:domain} fills the stored contact,
factor, product, and trace coordinates with these rational values. This proves completeness for every
integer input and every odd $B\ge3$.

\section{The polynomial and its exact certificate}

Each field value represents two rational coordinates. For the root abscissa
$a$, the four reciprocal squares and $C^2$, the stored second
coordinate is one half of the $\rho$ coefficient. Other stored field
values use their ordinary coefficients. These rational changes of
coordinates are invertible; they are fixed in the coefficient tables.
The six component equations of~\eqref{eq:H} are first made primitive,
then replaced by the nonsingular rational row combinations recorded in
the coefficient certificate. Their row-change determinant is
$-1/362963666515968$. Each resulting row, and each remaining component
equation, is cleared to a primitive integer quadratic, with positive
first nonzero coefficient in the prescribed monomial order.
These operations preserve their common zero set. Their sum of squares
defines $P_B$.

The coordinate count is
\begin{center}
\begin{tabular}{lr}
\toprule
Coordinates & Number\\\midrule
Input, its rational comparison component, and the common slope &3\\
Root chart for $A$ &6\\
Common chart for $H,H-P$ &8\\
Pair chart for $U,V_1$ &10\\
The two ratios &4\\
Four reciprocal squares and two comparison squares &12\\
Three scalar comparison-scale outputs &3\\
Two detector blocks, including their pole coordinates &42\\
The second argument $a_1\lambda$ &1\\\midrule
Total &89\\\bottomrule
\end{tabular}
\end{center}
Each detector uses three contact outputs $r_0,u_0,v_0$, three inverse
weights, four factor components, two product components, the two scalars
$\eta,b$, and seven pole coordinates: twenty-one in all. There are
$77$ primitive quadratic equations, so over $\Q$ their sum of squares
vanishes exactly when they all vanish.

The complete specification is in
\nolinkurl{computation/quartic/family_abc89_v21}. The producer and
independent checker are
\nolinkurl{build_abc_inverse_contact_v21.py} and
\nolinkurl{check_abc_inverse_contact_v21.py}. They specify all $77$ rows
and every coefficient of the full quartic. There are $937$ terms when
powers of $B$ are stored separately, and $936$ ordinary monomials for
every positive integer $B$. The fourth power of the rational component
of the root coordinate $h$ has coefficient one, proving degree four
and content one. Only the constant coefficient contains two powers of
$B$; it is $B^2+797681$. The remaining three parameter coefficients
are $4B,-2B,2B$, at the monomials $X^2$, $L_0^{\mathrm{prim}}$,
and $I_0$, respectively. Every other coefficient is independent of
$B$, and their largest absolute value is $4153456$. Thus
\[
 H(P_B)=\max(4153456,B^2+797681)
       =B^2+797681\quad\text{for odd }B\ge1833.
\]
No coefficient cancellation changes the stated support for positive
integer $B$. This coefficient threshold is distinct from the
ineffective sufficient defining shift in Theorem~\ref{thm:main}.

The independent checker imports no producer. It reconstructs all
24 changed contact, factor, product, and pole rows directly from the
displayed formulas, and compares the remaining 53 rows exactly with
their pinned coefficient source. It verifies every occurrence of the
omitted weight coordinates and independently expands all 77 quadratic
squares. Separate polynomial identities check the affine contact
translations, the contact normalization, the homogeneous trace scaling,
and preservation of the pole input. These identities hold before
assuming a product or trace parameter nonzero. The zero-weight and
zero-contact exclusions are the algebraic arguments in
Lemma~\ref{lem:domain}.

The common-point reconstruction dependency verifies the nonsingularity
of the rational row change. The separate standard-library program
\nolinkurl{computation/elliptic_constraints/check_fixed_difference_chart.py}
verifies twelve general polynomial identities, including elimination
without division, the two curve equations, translation by $P$, and the
exceptional leading-coefficient cases. These are finite algebraic
certificates. The arithmetic assumptions and the ineffective sufficient
shift remain those stated in Section~1.

\begin{thebibliography}{10}
\bibitem{Granville}
A.~Granville, Rational and integral points on quadratic twists of a given
hyperelliptic curve, \emph{International Mathematics Research Notices}
(2007), article rnm027, Proposition~1.
\url{https://doi.org/10.1093/imrn/rnm027}; author's manuscript,
\url{https://dms.umontreal.ca/~andrew/PDF/hyperelliptics.pdf}.
\bibitem{Dependencies}
\emph{Arithmetic dependencies of the construction with 29 variables},
companion paper, \nolinkurl{papers/arithmetic_dependencies}.
\bibitem{Selection}
\emph{Separation and denominator conditions in quadratic prime patterns},
companion paper, \nolinkurl{papers/prime_selection}.
\bibitem{Squares}
\emph{Square abscissas and direct charts for the elliptic subgroup},
companion paper, \nolinkurl{papers/square_abscissa}.
\bibitem{Trace}
\emph{Aggregate trace heights of elliptic multiples}, companion paper,
\nolinkurl{papers/trace_heights}; it applies J.~von Buhren,
\emph{J. Th\'eor. Nombres Bordeaux} \textbf{29} (2017), 307--320,
Theorem~1.1, \url{https://www.numdam.org/item/10.5802/jtnb.981.pdf}.
\bibitem{Logs}
\emph{Uniform formal-logarithm estimates}, companion paper,
\nolinkurl{papers/elliptic_logarithms}.
\bibitem{Source27}
\emph{A defining polynomial in twenty-seven variables: an existence
theorem with a uniform degree bound}, companion paper,
\nolinkurl{papers/few_variables27}, nonzero pole realization lemma.
\bibitem{PoleDensity}
\emph{Nonzero witnesses in the pole equations}, companion paper,
\nolinkurl{papers/pole_parameters/pole_witness_density_20261109_v14.tex}.
\end{thebibliography}
\end{document}
